\section{Structural and arithmetic boundaries}\label{sec:boundaries}
The preceding algorithms use more than a small leader dimension or a convex
follower. This section separates four restrictions: the interaction matrix in
the follower objective, the number and interaction graph of the leaders, the
coupling of follower constraints, and the requested arithmetic precision.
Throughout, \(\clip(t)=\min\{1,\max\{0,t\}\}\).

\subsection{A scalar leader with a dense positive definite follower}
\label{subsec:dense-hardness}
Consider the rational problem
\begin{equation}\label{eq:dense-box-model}
 \min_{x\in X} a\trans z(x)+b\trans x,\qquad
 z(x)=\argmin_{z\in[0,1]^N}
 \left\{\tfrac12z\trans Qz+(c+Dx)\trans z\right\},\quad Q\succ0.
\end{equation}
There are no upper constraints involving the response in this subsection.
The response is unique and continuous: any convergent sequence of leaders has
follower subsequences, and the limiting optimality inequalities identify every
subsequence limit with the unique response at the limiting leader. Thus a
compact leader polytope has an attained upper optimum.

Scalar-leader hardness and ternary digit gadgets have a close linear-bilevel
precedent in \citet[Theorem~2.1 and Section~4]{SugishitaCarvalho2026}. Exponential
strongly convex response paths also precede the following reduction.
In particular, the full-support sign patterns in
\citet[Proposition~2, equation~(4)]{MairalYu2012} imply a Boolean-vertex box path.
For their invertible square Lasso design \(A\) and target \(v\), the normalized
dual, with \(t=1/\lambda\), is
\[
 \min_{u\in[-1,1]^p}\tfrac12u\trans A^{-1}A^{-\mathsf T}u
       -t(A^{-1}v)\trans u.
\]
At a full-support primal sign pattern its dual point is that sign vector.
Their recursion visits every full sign vector whose last coordinate is positive;
an affine box change therefore gives every Boolean assignment in the first
\(p-1\) coordinates. This consequence of duality does not itself establish the
encoding bounds or upper-objective gap of a hardness reduction.

\begin{lemma}[Rational certificates for box responses]\label{lem:box-np}
For rational data in \eqref{eq:dense-box-model}, a rational upper threshold and a
rational leader polytope, threshold feasibility is in NP, even when the leader
dimension grows and affine upper rows are allowed.
\end{lemma}
\begin{proof}
Guess a partition \((L,F,U)\) of the coordinates into lower, free and upper
statuses. Set \(z_L=0,z_U=1\), and solve
\[
 z_F(x)=-Q_{FF}^{-1}(c_F+D_Fx+Q_{FU}\mathbf1).
\]
The principal matrix is positive definite; Cramer's rule gives affine rational
functions of polynomial bit length. Their box bounds, the endpoint gradient
signs \((Qz+c+Dx)_L\ge0\), \((Qz+c+Dx)_U\le0\), and all upper rows and the
threshold form a rational polyhedron in \(x\). If nonempty, it has a rational
point of polynomial bit length. Weak bounds on the free coordinates include
degeneracies. A verifier reconstructs \(z\) and checks these conditions exactly;
convexity makes them sufficient for global optimality. For a scalar leader the
polyhedron is simply an interval, possibly a singleton.
\end{proof}

\begin{theorem}[Dense scalar hardness with a constant gap]\label{thm:dense-hardness}
Threshold feasibility in \eqref{eq:dense-box-model} is NP-complete already for
\(X=[0,1]\), a rational positive definite \(Q\) independent of the leader, and an upper objective
with coefficients in \(\{-2,0,2\}\). A polynomial reduction gives an attained
optimum equal to zero in the yes case and at least two in the no case. All
leaders are feasible. The follower objective can be represented as a jointly
convex sum of affine squares in the leader and followers.
\end{theorem}
\begin{proof}
Start from a 3SAT instance with \(n\ge1\) variables and \(m\) clauses, each on
three distinct variables. This restriction preserves NP-hardness: remove
repetitions and tautologies; pad a two-literal clause using both signs of a
fresh variable, and a unit clause using all four sign pairs of two fresh
variables. An empty clause can be replaced by the fixed unsatisfiable formula
containing all eight clauses on three variables. Write \(\ell_a(y)\) for the
sum of the three continuous literal values in clause \(a\).
Put
\[
 \rho=\frac1{100\,3^n},\quad w_i=\rho^{i-1},\quad
 \eta=\rho^n,\quad \xi=\frac\eta{m+1},\qquad
 r_i=y_i-3^ix+2\sum_{j<i}3^{i-j}y_j+1.
\]
Use followers \((y,p,v)\in[0,1]^{2n+m}\) and minimize
\begin{equation}\label{eq:dense-square-gadget}
 \frac12\sum_iw_ir_i^2+
 \frac\eta2\sum_i(p_i-2y_i+1)^2+
 \frac\xi2\sum_a(v_a-1+\ell_a(y))^2.
\end{equation}
The residual map in the followers is square block triangular with diagonal
blocks \(L,I,I\), where \(L\) is lower triangular with unit diagonal.
Its positive weighted Gram matrix is positive definite.

Every Boolean vector \(d\in\{0,1\}^n\) occurs in the response at
\[
 x_d=\sum_{i=1}^n\frac{2d_i}{3^i}+\frac1{2\,3^n}\in(0,1).
\]
Indeed, with \(t_i=3^{i-1}x_d-2\sum_{j<i}3^{i-1-j}d_j\), the residual
\(r_i=d_i-3t_i+1\) has sign \(1-2d_i\) and magnitude between
\(1/(2\,3^{n-i})\) and one. The gradient of the first sum in coordinate \(i\)
is \(G_i=w_ir_i+\sum_{k>i}2\,3^{k-i}w_kr_k\). Its later-coordinate contribution
has magnitude at most
\[
 w_i\frac{6\rho}{1-3\rho}<\frac{w_i}{10\,3^n},
 \qquad (1-2d_i)G_i>\frac{w_i}{3^n}.
\]
At any complete follower optimum, conditional minimization gives
\begin{equation}\label{eq:shortfall-readout}
 p_i=\clip(2y_i-1),\quad y_i-p_i=\min(y_i,1-y_i),\quad
 v_a=\max(0,1-\ell_a(y)).
\end{equation}
At \(y=d\), take these conditional optima. The \(p\)-squares add derivative
\(-2\eta\) at a zero bit and zero at a one bit. A clause residual is respectively
\(0,0,1,2\) when the clause has \(0,1,2,3\) true literals, so the magnitude of
all clause feedback at any base coordinate is at most \(2m\xi<2\eta\).
The total perturbation is less than
\(4\eta=w_n/(25\,3^n)\), below the base margin. All box optimality signs
therefore hold, including those of the auxiliary coordinates. Positive
definiteness proves that this is the unique global response.

Choose the upper objective
\[
 H=2\sum_i y_i-2\sum_i p_i+2\sum_a v_a.
\]
It is nonnegative by \eqref{eq:shortfall-readout}, and a satisfying Boolean
response gives zero. If the formula is unsatisfiable, round \(y\) at \(1/2\)
and take a false clause. Its literal sum is at most
\(D_y=\sum_i\min(y_i,1-y_i)\), since its variables are distinct. Consequently
\[
 H\ge 2D_y+2\max(0,1-D_y)\ge2.
\]
The shortfall coordinates have thus replaced all upper clause constraints by
an objective penalty inside the response. No upper feasibility assumption is
needed. Every coefficient and every displayed witness has polynomial bit
length: powers \(3^i\) have \(O(n)\) bits and the weights have \(O(n^2)\) bits.
Expanding the squares yields \eqref{eq:dense-box-model} after removing a
leader-only term. The unexpanded squares are jointly convex; deleting that
term preserves responses but need not preserve joint convexity in \((x,z)\).
Lemma~\ref{lem:box-np} proves NP membership.
\end{proof}

\subsection{Exact hardness arbitrarily close to the identity}
\label{subsec:conditioned-hardness}
The small weights in the last proof do not control conditioning. A different
construction does, at the cost of a much smaller upper gap. Positive diagonal
network scaling is established in \citet[equations~(2.6)--(2.7)]{ElGhaouiEtAl2021},
and convex squared-residual energies with weak backward feedback are used by
\citet[Section~2]{ZachEstellers2019}. The proof below supplies explicit rational
scales and a relative-coordinate error bound needed for the reduction.

\begin{theorem}[Near-identity exact hardness]\label{thm:conditioned-hardness}
Threshold feasibility in \eqref{eq:dense-box-model} with \(X=[0,1]\) remains NP-complete when
all normalized follower and upper coefficients have magnitude at most two and
\[
 (49/50)^2 I\preceq Q\preceq(51/50)^2 I,\qquad \kappa_2(Q)<2.
\]
For any supplied rational \(0<\zeta\le1\), the reduction can also ensure
\(\norm{Q-I}_1,\norm{Q-I}_\infty,\norm{Q-I}_2<\zeta\) and strict diagonal
dominance. The resulting yes/no objective gap is exponentially small but has
polynomial bit length. Unless \(\mathrm P=\mathrm{NP}\), there is no global
additive approximation algorithm polynomial in input length and accuracy bit
length for this class.
\end{theorem}
\begin{proof}
Use the same clause functions. Starting with \(t_1=x\), define a feedforward
ReLU computation
\[
 \begin{gathered}
 \alpha_i=(3t_i-1)_+,\quad \beta_i=(3t_i-2)_+,\quad y_i=\alpha_i-\beta_i,\\
 t_{i+1}=3t_i-2y_i,\qquad p_i=(2y_i-1)_+,\qquad v_a=(1-\ell_a(y))_+.
 \end{gathered}
\]
The map \(t\mapsto3t-2\clip(3t-1)\) sends \([0,1]\) into itself.
Order the states as \(h=(\alpha_1,\beta_1,\ldots,\alpha_n,\beta_n,p,v)\),
with \(N=3n+m\). Substituting the recursion for each \(t_i\) gives
\[
 h=(Ah+b_0+xb_1)_+,
\]
where \(A\) is strictly lower triangular and its entries and those of \(b_0,b_1\)
have magnitude at most \(C=2\,3^n\). On \([0,1]\),
\(0\le h\le2\mathbf1\), and
\(0\le h-Ah-b_0-xb_1\le2\mathbf1\): the preactivations for \(\alpha,\beta,p,v\)
lie respectively in \([-1,2],[-2,1],[-1,1],[-2,1]\).
The score
\[
 \Phi(h)=2\sum_i\alpha_i-2\sum_i\beta_i-2\sum_i p_i+2\sum_a v_a
          =\ell\trans h
\]
has \(\norm\ell_1=2N\). The previous rounding argument gives \(\Phi\ge0\),
value zero at a satisfying encoded leader, and \(\Phi\ge2\) everywhere for an
unsatisfiable formula. At \(x_d\), the ternary recursion gives \(y=d\).

Set
\[
 M=(1+NC)^N,\quad \theta=\frac1{100NCM},\quad
 s_i=\frac{\theta^{i-1}}{4C},\quad S=\diag(s_i),\quad B=SAS^{-1}.
\]
If the additional tolerance is requested, replace \(\theta\) by
\(\min\{\theta,\zeta/(10C)\}\). The unit-box follower minimizes
\begin{equation}\label{eq:scaled-network-energy}
 \tfrac12\norm{(I-B)u-S(b_0+xb_1)}_2^2.
\end{equation}
Since \(|B_{ij}|\le C\theta^{i-j}\) for \(i>j\), both its row and column norm
are at most \(2C\theta\le1/50\). Thus \(Q=(I-B)\trans(I-B)\) has the asserted
spectral bounds. Moreover \(\norm{Sb_j}_\infty\le1/4\), so the normalized
linear-cost coefficients are less than one in magnitude. The spectral bound
also bounds every entry of \(Q\) by two.

Here the small absolute size of \(u\) is not an adequate error measure.
Let \(u\) be the true minimizer, \(e_i=|u_i/s_i-h_i|\),
\(P=|A|\), and \(r=(I-B)u-S(b_0+xb_1)\). The exact KKT fixed point is
\[
 u=\clip\bigl(Bu+S(b_0+xb_1)+B\trans r\bigr).
\]
The scaled feedforward point \(Sh\) is the same clipped fixed point without
\(B\trans r\); its positive coordinates are strictly below one.
Coordinatewise nonexpansiveness of clipping gives
\[
 e\le Pe+U|S^{-1}r|,\qquad
 U=S^{-1}|B|\trans S=S^{-2}P\trans S^2,
 \qquad |S^{-1}r|\le(I+P)e+2\mathbf1.
\]
The nonnegative matrix \(T=(I-P)^{-1}=\sum_{j=0}^{N-1}P^j\) has
\(\norm T_\infty\le M\); also \(\norm P_\infty\le NC\) and
\(\norm U_\infty\le2C\theta^2\). Multiplication by \(T\) therefore yields
\[
 \norm e_\infty\le 2MC(1+NC)\theta^2\norm e_\infty+4MC\theta^2,
 \quad
 \norm e_\infty\le8MC\theta^2\le\frac1{16N},
\]
since the first coefficient is less than \(1/2\). These inequalities continue
to hold after decreasing \(\theta\).
Put \(\delta=s_N\) and use upper objective \(G(u)=\delta\ell\trans S^{-1}u\).
Every coefficient has magnitude at most two, while
\[
 |G(u)-\delta\Phi(h)|\le\delta/8.
\]
The yes optimum is at most \(\delta/8\); the no optimum is at least
\(15\delta/8\). Comparing with threshold \(\delta\) proves hardness, and
Lemma~\ref{lem:box-np} supplies membership. All scales have polynomial bit
length, including the additional input length of \(\zeta\).
Finally, \(\norm B_1,\norm B_\infty\le\zeta/5\) gives
\(\norm{Q-I}_{1,\infty,2}\le2\zeta/5+\zeta^2/25<\zeta\).
For each row, \(Q_{ii}-\sum_{j\ne i}|Q_{ij}|\ge1-\norm{Q-I}_\infty>0\).
Accuracy \(\delta/4\) separates the two cases and has polynomially many bits.
\end{proof}
The last theorem is ordinary exact hardness. Its bounded coefficients and
conditioning do not yield strong NP-hardness or an inverse-polynomial
absolute gap under the same restrictions. In particular it is compatible with
the following algorithm, whose dependence is on inverse accuracy.

\subsection{A conditioned approximation scheme for dense box followers}
\label{subsec:conditioned-algorithm}
\begin{theorem}[A cost-space grid]\label{thm:conditioned-grid}
Fix \(r\). For \eqref{eq:dense-box-model}, let \(X\subseteq[0,1]^r\) be a
nonempty rational polytope. For rational \(0<\eps\le1\), one can return a
rational leader and its exact rational follower with upper value at most
\(\mathrm{OPT}+\eps\norm a_1\), in time polynomial in the rational input length,
\(K=\norm Q_\infty\norm{Q^{-1}}_\infty\), and \(1/\eps\). Here
\(K\le N\kappa_2(Q)\) for \(N\ge1\). There are no response-dependent upper rows.
\end{theorem}
The maximum-volume basis used below is the classical barycentric-spanner
argument of \citet[Proposition~2.2]{AwerbuchKleinberg2004}. Approximation of
parametric convex solution paths and uniform optimizer errors have earlier
treatments, including \citet{GiesenJaggiLaue2012} and
\citet[Theorem~4]{BemporadFilippi2001}. The particular purpose of the saturation
step here is to remove numerical dependence on the affine-cost magnitudes;
all those coefficients still contribute their binary encoding lengths.
The scalar additive path bound of \citet[Lemma~4 and Theorem~5]{GiesenEtAl2012}
is another close precedent: its coefficient involves parameter range and slope
variation of the follower value. Strong convexity converts follower-objective
error into response error, but alone does not remove those numerical factors.

\begin{proof}
If \(N=0\), solve the leader LP. If \(r=0\), evaluate the sole follower. If
\(a=0\), solve the leader LP and evaluate its follower for the promised output.
Hence assume \(N,r\ge1\) and \(\norm a_1>0\). Put
\(\mu=1/\norm{Q^{-1}}_\infty\le\lambda_{\min}(Q)\). For row \(i\), write
\[
 m_i^- =\sum_j\min(Q_{ij},0),\quad m_i^+=\sum_j\max(Q_{ij},0),\quad
 R_i=m_i^+-m_i^-=\sum_j|Q_{ij}|>0.
\]
If its linear cost \(t_i\ge-m_i^-\), its response coordinate is zero:
for \(z_i>0\) its gradient is at least \(Q_{ii}z_i+t_i+m_i^->0\),
contradicting box optimality. Similarly \(t_i\le-m_i^+\) forces \(z_i=1\),
because at \(z_i<1\) the gradient is at most
\(Q_{ii}(z_i-1)+t_i+m_i^+<0\). These tests include equality at thresholds.

Arrange the \(2N\) hyperplanes \(c_i+D_ix=-m_i^\pm\), intersect all relative
cells with \(X\), and use their closures. There are \(N^{O(r)}\) such cells.
On each, coordinates outside a set \(J\) are fixed at endpoints; all costs
with indices in \(J\) lie in the closed slabs of width \(R_i\). Coordinates
with constant costs inside a slab remain in \(J\). For two leaders in this
cell, the two variational inequalities and positive definiteness give
\begin{equation}\label{eq:cost-grid-sensitivity}
 \norm{z(x)-z(x')}_2\le\frac{\norm{D_J(x-x')}_2}{\mu}.
\end{equation}
Indeed the response difference vanishes outside \(J\), so the cost difference
in those coordinates contributes nothing to their summed inequalities.

Let \(E=D_J/\mu\) have rank \(q\le r\). If \(q=0\), the response is constant
on the cell and minimizing \(b\trans x\) suffices. Otherwise choose \(q\)
independent columns of \(E\), then a maximum absolute \(q\)-row minor in those
columns. Let \(B_0\) denote the selected full rows. Every row of \(E\) is a
linear combination of those rows with coefficients in \([-1,1]\): replacing
one row in the maximal minor and applying Cramer's rule proves the bound;
the selected columns determine dependencies among full rows. Consequently
\[
 \norm{B_0(x-x')}_\infty\le h
 \quad\Longrightarrow\quad \norm{z(x)-z(x')}_2\le\sqrt N\,qh.
\]
A selected coordinate \(D_ix/\mu\) ranges in
\([(-m_i^+-c_i)/\mu,(-m_i^--c_i)/\mu]\), of width
\(R_i/\mu\le K\). Its offset may be large but has polynomial bit length.
Divide each of these \(q\) intervals into closed intervals of length at most
\(h=\eps/(Nr)\). For every resulting rectangle, intersect its inverse image
under \(B_0\) with the current cell, minimize \(b\trans x\) there by rational
LP, and evaluate the exact follower. Keep the best value. The number of LPs is
\(N^{O(r)}(2+KNr/\eps)^r\), and their data have polynomial bit length.
An optimal leader belongs to one of these closed sets. Its LP representative
has no larger direct leader cost and changes the response term by at most
\(\norm a_1\sqrt N\,rh\le\eps\norm a_1\). Cell boundaries are included by the
weak saturation tests.

For completeness an elementary exact follower oracle fits the claimed
\(\poly(K)\) running time. At a rational leader put \(t=c+Dx\) and
\(L_Q=\norm Q_\infty\). Projected gradient iteration from zero,
\[
 z^{j+1}=\clip\bigl(z^j-(Qz^j+t)/L_Q\bigr),
\]
contracts Euclidean distance to the optimum by at most \(1-1/K\). This follows
from nonexpansiveness of projection and the eigenvalues of \(I-Q/L_Q\).
Clear all denominators in \(Q,t\) simultaneously, and bound the absolute
integer entries by \(H\ge1\). In any active principal system, the denominator
of each optimum coordinate in reduced form is at most \(B=N!H^N\);
Cramer's rule gives this bound also when free coordinates reach endpoints.
After
\[
 j=\left\lceil K\bigl(2\lceil\log_2B\rceil+
               \lceil\log_2N\rceil+4\bigr)\right\rceil
\]
iterations every coordinate is within \(1/(4B^2)\) of its optimum. Two distinct
rationals of denominator at most \(B\) differ by at least \(1/B^2\), so the
optimum is the unique such rational in that error interval. It is found among
the continued-fraction convergents of the iterate: the error is less than
\(1/(2q^2)\) for its denominator \(q\le B\); endpoints zero and one can be
checked separately. Verify the recovered vector by exact box KKT conditions.
The iteration count is polynomial in \(K\) and input length; rational iterate
bit lengths grow at most linearly in that count times the input length.
All recovery and verification therefore have the claimed bit complexity.
Finally, symmetry gives \(\norm Q_\infty\le\sqrt N\norm Q_2\) and the same
inequality for \(Q^{-1}\), establishing \(K\le N\kappa_2(Q)\).
\end{proof}
A grid directly in leader coordinates would not have this uniform guarantee.
For \(Q=2^{-L}\) and linear cost \(-xz\), the scalar response is
\(\clip(2^Lx)\), although its Hessian condition number is one. Saturation and
cost normalization isolate precisely this narrow transition. Neither the
number of grid cells nor the runtime above is polynomial in accuracy bit
length, or fixed-parameter tractable in \(r\).

\subsection{Growing leader dimension and a bounded core}
\label{subsec:leader-structure}
With an identity Hessian and independent boxes, each follower is
\(z_i(x)=\clip(h_i(x))\) for an affine rational \(h_i\). Thus affine upper
optimization on a leader product box is optimization of a signed sum of
clipped affine functions. Define its \emph{leader interaction graph} by joining
two leader coordinates whenever both occur in some \(h_i\). This graph refers
to cost dependence on leaders, not to constraints joining follower coordinates.

\begin{theorem}[Growing leader dimension]\label{thm:growing-leader}
With growing leader dimension \(r\), identity follower Hessian, independent
unit boxes and no upper constraints, threshold optimization is NP-complete
and W[1]-hard parameterized by \(r\). Under the exponential time hypothesis
(ETH), no algorithm has runtime \(f(r)\mathcal L^{o(r)}\), where \(\mathcal L\)
is input length. These conclusions hold with each follower depending on at
most two leaders, all follower affine coefficients of magnitude at most one,
and upper coefficients of magnitude at most two. The reduction has a constant
absolute upper-objective gap.
\end{theorem}
\begin{proof}
The parameterized source result is due to
\citet[Proposition~4.1, Theorem~5.3 and Corollary~5.5]{FroeseEtAl2026}.
From Multicolored Clique with \(k\) colors it constructs a two-layer ReLU
network in \(k\) inputs with maximum \(C_0=k+\binom k2\) in the yes case and
at most \(C_0-1\) otherwise. Its realizing points lie in a polynomially bounded
box \([0,L]^k\); all affine coefficients are polynomially bounded in magnitude,
and each neuron depends on at most two inputs. Output coefficients can be
made \(\pm1\) by positive homogeneity and duplication of their bounded
integer magnitudes. These properties follow directly from its node and edge
spikes, whose locations are integers in a greedy Sidon set of polynomial
magnitude. We use the proved gap, not an assumption that continuous inputs
must themselves be discrete.

Rescale the input as \(t=Lx\). Write the resulting network as
\(\sum_i\sigma_i(h_i(x))_+\), with \(\sigma_i\in\{-1,1\}\). Choose a positive
integer \(R>\max_i(|h_i(0)|+\norm{\nabla h_i}_1)\), polynomially bounded in
the source size. Then
\[
 \clip(h_i(x)/R)=(h_i(x))_+/R\quad(x\in[0,1]^k).
\]
Each coordinate is the response of the identity-Hessian follower with cost
\(-h_i(x)z_i/R\). The upper expression
\(C_0-R\sum_i\sigma_i z_i(x)\) has gap zero versus at least one.
Duplicate each coordinate \(R\) times to replace its coefficient \(R\sigma_i\)
by unit coefficients. If a constant-free linear objective is required, add
\(C_0\) fixed-one responses. All follower coefficients now have magnitude at
most one; the duplication is polynomial and preserves dependence on at most
two leaders. The source parameter and ETH bounds transfer with \(r=k\).
Lemma~\ref{lem:box-np} proves ordinary NP membership.
\end{proof}
This theorem transfers an established parameterized hardness result. An
alternative explicit construction below records a useful continuous rounding
gap and avoids reliance on the particular shape of the source spikes.

\begin{proposition}[A mixed-radix gap]\label{prop:mixed-radix}
Given a graph on \(n\ge2\) labeled vertices and \(k\ge2\), one can construct in
polynomial time a nonnegative signed clipped-affine objective on \([0,1]^k\)
whose minimum is zero for a \(k\)-clique and at least two otherwise. It has an
identity-Hessian realization satisfying the coefficient and two-leader
restrictions of Theorem~\ref{thm:growing-leader}.
\end{proposition}
\begin{proof}
Put \(t_i=(n-1)x_i\), and let \(d(t)\) be distance to the nearest integer label
in \(\{0,\ldots,n-1\}\). On this interval,
\[
 2d(t)=\sum_{a=0}^{2n-3}(-1)^a\clip(2t-a).
\]
For integer \(s=u+nv\), define \(f(s)=0\) if \(u,v\) are distinct adjacent
vertices and \(f(s)=1\) otherwise. Its piecewise linear interpolation is
\[
 \varphi(s)=1+\sum_{a=0}^{n^2-2}(f(a+1)-f(a))\clip(s-a),
 \quad 0\le s\le n^2-1.
\]
It lies in \([0,1]\) and is 1-Lipschitz. Take
\[
 H(t)=2n\sum_i d(t_i)+2\sum_{i<j}\varphi(t_i+nt_j).
\]
A clique labeling gives zero. Otherwise round every \(t_i\) to a label, let
\(D_t=\sum_i d(t_i)\), and choose a bad pair of rounded labels. For that pair
\(\varphi(t_i+nt_j)\ge\max(0,1-nD_t)\). Hence
\(H\ge2nD_t+2\max(0,1-nD_t)\ge2\).
The displayed formulas use \(O(k^2n^2)\) clipped coordinates after duplicating
the coefficient \(n\) in the first sum; upper coefficients have magnitude at
most two. Every affine argument depends on at most two leaders and has
polynomially bounded integer coefficients. To bound those coefficients by one,
choose \(R=2n^2\) and use, for each such argument \(h\),
\[
 \clip(h)=R\{\clip(h/R)-\clip((h-1)/R)\}.
\]
On the leader box neither positive scaled argument exceeds one, so the
identity follows from \(\clip(h)=(h)_+-(h-1)_+\). Duplicate \(R\) copies to
retain upper coefficients of magnitude at most two. This gives
\(O(k^2n^4)\) coordinates with follower affine coefficients bounded by one.
Constant terms are realized by fixed-one responses if needed.
\end{proof}

\begin{theorem}[A supplied bounded leader core]\label{thm:leader-core}
Suppose the leader domain is a product of rational intervals, the follower
Hessian is positive diagonal, and all other data are as in
\eqref{eq:dense-box-model}, with no response-dependent upper rows. Suppose a
set of at most \(c\) leader vertices is supplied whose deletion leaves
components of size at most \(h\) in the leader interaction graph. For fixed
\(c,h\), exact optimization and a rational optimal leader-response pair can
be computed in polynomial time.
\end{theorem}
\begin{proof}
Normalize the leader intervals to unit intervals; fixed coordinates are removed.
Normalize each positive diagonal quadratic to obtain its clipped affine
response. Write the core as \(u\), and the component leaders as \(v_b\) with
\(h_b\le h\). The upper objective decomposes as
\(H_0(u)+\sum_bH_b(u,v_b)\), since a follower's noncore arguments all belong
to one component.
For fixed \(u\), an optimum of \(H_b\) on its box occurs at a vertex of the
arrangement of its clipping hyperplanes and box facets: it is affine on every
closed arrangement cell. Enumerate every choice of \(h_b\) hyperplanes with
independent coefficients in \(v_b\). Solving them gives a candidate
\(v_b=P u+q\), with rational affine data of polynomial bit length.
There are at most \((2m_b+2h_b)^{h_b}\) candidates, where \(m_b\) is the
number of clipped terms of that component. Discard singular choices. Box
facets ensure that all relevant arrangement vertices are included even when
some clipping hyperplanes coincide or miss the box at particular parameters.
A candidate is feasible precisely when it belongs to the component box; it
need not remain a vertex away from the parameter where it was generated.

In core space arrange the hyperplanes where a candidate coordinate reaches
zero or one, where a clipped term evaluated at a candidate reaches either
threshold, and where a core-only term reaches a threshold. On each relative
cell, candidate feasibility is constant and every feasible candidate's value is
a fixed affine function of \(u\). Add all pairwise equality hyperplanes between
candidate values for the same component, and refine the cell. A minimizing
candidate for each component is now fixed, and the summed objective is affine.
Minimize it by LP on the closure of each resulting cell and reconstruct the
selected candidates and followers. All enumerations have polynomial size for
fixed \(c,h\); constant core dimension also covers \(c=0\).

The closure step is sound: each selected candidate remains box-feasible and
its continuous clipped value agrees with its affine formula on the closure.
At a boundary, a newly feasible candidate omitted from this particular LP
could be better, but that does not invalidate the returned feasible value.
Completeness follows by considering the relative cell containing an optimal
core point, where all its feasible candidates are considered. Thus the best
constructed point is globally optimal. Rational LP recovery proves the output
claim.
\end{proof}
The vertex integrity of a graph is the minimum, over deleted vertex sets, of
their size plus the largest remaining component size. Thus the supplied core
in Theorem~\ref{thm:leader-core} certifies vertex integrity at most \(c+h\).
This is a bounded-core decomposition in the leader cost graph. Related
``few global variables'' decompositions for integer programming are studied
by \citet{DvorakEtAl2021}; their integer constraint graphs and conclusions are
different. The algorithm here is polynomial for fixed \(c,h\), with exponents
depending on these parameters. A core need not be found by an FPT algorithm:
enumerating all sets of size at most \(c\) also gives a polynomial algorithm
when \(c\) is fixed. If each follower depends on only one leader, the objective
separates completely and each univariate clipped sum can be minimized at its
breakpoints.

\begin{theorem}[A leader path is already hard]\label{thm:leader-path}
With a growing number of leaders, the identity-Hessian independent-box model
is weakly NP-hard even when its leader interaction graph is a path and all
follower affine coefficients have magnitude at most one and upper coefficients
have magnitude at most two. Threshold feasibility is NP-complete. Exact
continuous-state messages on a path can have exponentially many affine pieces,
even for an identical local factor with a fixed finite coefficient alphabet.
\end{theorem}
\begin{proof}
For positive Subset Sum weights \(a_1,\ldots,a_n\), let \(W=\sum_i a_i\).
Targets outside \([0,W]\), and its endpoints, can be handled directly; assume
\(0<T<W\). Put \(r_i=a_i/W\), \(\tau=T/W\). For \(-1\le t\le1\), define
\[
 \rho_r(t)=\min\{|t|,|t-r|\}
 =-t+2\clip(t)-2\clip(t-r/2)+2\clip(t-r),\quad0<r\le1.
\]
On leaders \(x_0,\ldots,x_n\in[0,1]\) take
\[
 H(x)=x_0+|x_n-\tau|+\sum_{i=1}^n\rho_{r_i}(x_i-x_{i-1}).
\]
Choose \(\delta_i\in\{0,r_i\}\) nearest to the increment and telescope to get
\(\bigl|\sum_i\delta_i-\tau\bigr|\le H(x)\). Conversely, cumulative sums of
any chosen subset lie in the box and give that subset's distance to \(\tau\).
Therefore
\begin{equation}\label{eq:subset-path-distance}
 \min H=\frac1W\min_{S\subseteq[n]}\left|\sum_{i\in S}a_i-T\right|.
\end{equation}
The clipped expansion is
\[
 2x_0-2x_n+\tau+2\clip(x_n-\tau)+
 2\sum_i\{\clip(x_i-x_{i-1})-\clip(x_i-x_{i-1}-r_i/2)
                         +\clip(x_i-x_{i-1}-r_i)\}.
\]
It uses \(3n+1\) independent identity-Hessian followers, each depending only
on consecutive leaders or the terminal leader. Unary copies of \(x_0,x_n\)
and a fixed-one response may replace the direct affine terms. The coefficient
bounds follow, and Lemma~\ref{lem:box-np} gives NP membership. The gap is
\(0\) versus at least \(1/W\). Scaling the objective by \(W\) produces a
constant gap but loses the bounded coefficient claim; no strong hardness is
inferred.

There is also an exact message identity. Let \(S_n\) be the set of normalized
subset sums, and minimize \(x_0+\sum_i\rho_{r_i}(x_i-x_{i-1})\) with \(x_n=t\).
The same telescope is a lower bound of \(\dist(t,S_n)\). For the reverse bound,
follow a zero-cost subset path to its terminal sum and change only its last
coordinate to \(t\). Its last factor increases by at most that distance.
The message is therefore exactly \(\dist(t,S_n)\). Weights \(a_i=2^{i-1}\)
give \(S_n=\{j/(2^n-1):0\le j\le2^n-1\}\), hence \(2(2^n-1)\) maximal
affine pieces on \([0,1]\).
For a fixed alphabet use instead the identical factor
\(\rho_{1/2}(x_i-x_{i-1}/2)\). Its zero-cost terminal states are
\(S_n=\{j/2^n:0\le j<2^n\}\). If \(e_i=x_i-x_{i-1}/2-\delta_i\),
\(\delta_i\in\{0,1/2\}\), then
\[
 t-\sum_i2^{-(n-i)}\delta_i
 =2^{-n}x_0+\sum_i2^{-(n-i)}e_i.
\]
The absolute value is bounded by \(x_0+\sum_i|e_i|\); the same last-coordinate
construction gives equality with \(\dist(t,S_n)\). There are
\(2^{n+1}-1\) maximal affine pieces, including the last tail to one.
All local coefficients now belong to one finite set (gains \(1,-1/2\), shifts
\(0,-1/4,-1/2\), and readout coefficients \(\pm2\)).
\end{proof}
Closure of continuous piecewise linear functions under variable elimination
is classical; see \citet[Lemmas~1--2 and Theorem~1]{MeuleauEtAl2008} for a
scheduling formulation. Closure alone gives no polynomial bound on intermediate
message size. The identical-factor family above has the trivial optimum zero;
its exponential message size is an output-complexity example, not another
NP-hardness proof.

\subsection{A path in the follower constraints}
\label{subsec:follower-path}
The preceding path connected \emph{leaders} in affine follower costs. Now there
is one leader, the Hessian is diagonal, and a path connects \emph{follower}
coordinates through inequalities. This distinction explains why the next
example lies outside the independent-block response theorem.
For rational \(0<\gamma<1/2\), set \(z_0=0\) and
\begin{equation}\label{eq:klee-polytope}
 P_n(\gamma)=\{z:0\le z_1\le1,\quad
          \gamma z_{j-1}\le z_j\le1-\gamma z_{j-1}\ (2\le j\le n)\}.
\end{equation}
Every row involves at most two consecutive coordinates (2VPI).
Appendix~\ref{app:path-geometry} proves the needed geometric and arithmetic
facts, and states with its proof the sparse-shadow construction of
\citet[Section~4]{GartnerEtAl2013}.

\begin{proposition}[What a sparse shadow obstructs]\label{prop:path-shadow-boundary}
For \(\gamma=1/4\), there is an explicit scalar-price LP over \(P_n(\gamma)\)
with \(2^n\) distinct optimal vertices, each uniquely optimal on a price
interval. Its value has \(2^n\) affine pieces, although evaluation at a
rational price requires only \(O(n)\) rational arithmetic operations of
polynomial bit length. The terminal-state value message has \(2^n-1\) affine
pieces. An unextended quantifier-free polynomial description of the bounded
epigraph in the two price/value variables requires the sum of the degrees of
its nonzero polynomials to be at least \(2^n\).
\end{proposition}
The explicit formulas and the boundary-factor proof are in
Appendix~\ref{app:path-geometry}. This is a limitation on enumerating an entire
response or on eliminating all auxiliary variables, not on solving one LP,
nor on arithmetic circuits or descriptions retaining quantified states.

\begin{theorem}[Scalar-leader hardness with an identity-Hessian path follower]
\label{thm:follower-path-hardness}
There is an NP-complete restricted family of rational bilevel feasibility
problems with one leader in \([-1,1]\), a unique follower minimizing a strictly
convex diagonal quadratic over a 2VPI path polytope, one upper slab with two
linear bounds, and one \emph{nonconvex quadratic upper row}. The follower
Hessian can be exactly the identity. Without upper rows, such path followers
can have at least \(2^n\) distinct constant response intervals.
The feasibility hardness also holds when all \emph{local follower inequality}
coefficients and right-hand sides belong to fixed finite alphabets.
\end{theorem}
\begin{proof}
Take Subset Sum data \(a_i>0\), \(W=\sum_i a_i\), \(0\le T\le W\), and put
\(\gamma=1/(8W)\). Define
\[
 c_i=(1-\gamma)\gamma^{2(n-i)-1}\ (i<n),\quad c_n=0,\quad
 F(z)=z_n-c\trans z-z_n^2.
\]
Lemma~\ref{lem:klee-telescope} shows that \(F\ge0\) on \(P_n(\gamma)\), with
zeros precisely at its vertices. Each vertex has a bit description
\(v_j=d_j+(1-2d_j)\gamma v_{j-1}\), and
\(|v_j-d_j|\le\gamma\). The slab
\begin{equation}\label{eq:klee-slab}
 T-1/4\le\sum_i a_i z_i\le T+1/4
\end{equation}
therefore contains a vertex if and only if the Subset Sum instance is positive:
the forward error is at most \(\gamma W=1/8\), and conversely a vertex in
the slab gives \(|\sum_i a_i d_i-T|\le3/8<1\), forcing equality of integers.

It remains to make every required vertex an exact strongly convex response.
Terminal vertex values are distinct and have denominators dividing
\(D=(8W)^{n-1}\). Thus their separation is at least \(1/D\). Put
\(\Delta=D^{-2}\), \(\tau=\Delta/(2n)\), and let the follower minimize
\begin{equation}\label{eq:klee-quadratic-follower}
 \tfrac\tau2\norm z_2^2-c\trans z-\lambda z_n
 \quad\text{over }P_n(\gamma),\qquad -1\le\lambda\le1.
\end{equation}
At vertex \(v\) choose \(\lambda_v=2v_n-1\). Since
\(c\trans u=u_n-u_n^2\) at every vertex \(u\), the loss in the linear score
from \(v\) to any other vertex is \((u_n-v_n)^2\ge\Delta\).
For an arbitrary convex combination \(z=\sum_u\alpha_u u\), put
\(a=\sum_{u\ne v}\alpha_u\). Then the score loss is at least \(\Delta a\),
\(\norm{z-v}_\infty\le a\), and
\(\norm z_2^2-\norm v_2^2\ge-2na\). Consequently its follower objective
exceeds that at \(v\) by at least \((\Delta-\tau n)a=\Delta a/2\), strictly
if \(z\ne v\). The same reasoning at price \(\lambda_v+h\) loses at most
\(|h|a\), so \(v\) remains the unique response whenever \(|h|\le\Delta/4\)
within the leader interval. There is a nontrivial one-sided interval even at
an endpoint. Fixing \(\gamma=1/4\) proves the exponential constant-response
claim with no upper constraints.

Impose \(F(z)\le0\) and \eqref{eq:klee-slab} as upper rows. The former forces
an original path vertex; the preceding equivalence and exposing leaders give
the reduction in both directions. Guessing its bits, computing the vertex and
\(\lambda_v\), and checking the slab is a polynomial-bit NP certificate for
this restricted family. All coefficients have polynomial bit length, including
\(1/\tau\). Multiplication of the entire follower objective by \(1/\tau\)
preserves the response and makes its Hessian the identity.

For the fixed local alphabet, apply the padding construction of
Lemma~\ref{lem:klee-padding}: use \(\gamma=1/4\), dimension
\(N=n+(n-1)r\) with \(4^{-(r+1)}\le1/(8W)\), and bound each padding
coordinate by \(1/2\). Use the full \(N\)-dimensional cube's \(F,c\), and
\(D=4^{N-1},\Delta=D^{-2},\tau=\Delta/(2N)\). Every required free-bit
pattern is an original cube vertex satisfying the padding bounds. Its strict
response inequality was proved over the entire cube, so it remains valid on
the padded subset. Conversely \(F\le0\) forces an original cube vertex and
its padding bits must be zero. Apply the slab only to the free coordinates;
the same \(1/8\) error proves the equivalence. Local inequality coefficients
are now in \(\{0,\pm1,\pm1/4\}\), and right-hand sides in \(\{0,1/2,1\}\).
\end{proof}
The upper quadratic row has Hessian \(-2e_ne_n\trans\); it is not a convex
upper constraint. The same reduction can minimize \(F(z(\lambda))\) under the
linear slab at threshold zero. The slab-only feasible set is nonempty: the
continuous response visits zero and the all-one-bit vertex, whose weighted
free-coordinate sum is at least \(W-1/8\); the intermediate value theorem
covers integer targets below \(W\), and that last vertex covers \(T=W\).
Its compact feasible graph gives attainment. Deleting the quadratic row from
the feasibility version therefore makes every target feasible. This proof
establishes neither all-linear-upper feasibility hardness nor strong hardness.
The fixed local alphabet does not bound the slab weights or the rescaled
linear costs. Rank-one concave quadratic hardness more generally has the
classical precedent \citet{PardalosVavasis1991}.

There is a narrow positive counterpart. If both sides of each affine interval
transition use the \emph{same} gain, exact feasibility projection along paths
is polynomial, including negative or zero gains. Proposition~\ref{prop:strip-projection}
in Appendix~\ref{app:strip-projection} gives the full construction, its extension
to trees and fixed \emph{total} cycle rank, and same-field witness recovery.
The opposite gains \(+\gamma\) and \(-\gamma\) in
\eqref{eq:klee-polytope} do not satisfy that hypothesis. Nor does that
projection permit arbitrary objectives involving all eliminated states.

\subsection{Curvature, degree encoding, and arithmetic output}
\label{subsec:arithmetic-boundaries}
\begin{proposition}[Removing positive local curvature]\label{prop:curvature-boundary}
Optimistic feasibility is NP-complete in restricted degree-two families with
no leader and an independent unit-box follower if either zero local curvature
and quadratic upper rows, or negative local curvature and linear upper rows,
are permitted.
\end{proposition}
\begin{proof}
With zero follower objective, every \(z\in[0,1]^N\) is optimal. Upper rows
\(z_i(1-z_i)=0\) and the linear literal-sum inequalities \(\ell_a(z)\ge1\)
encode 3SAT. Alternatively minimize \(\sum_i z_i(1-z_i)\) in the follower.
Its minimum is zero and its minimizers are exactly the Boolean vertices;
then only the linear upper clause rows are needed. Boolean assignments are
certificates for both restricted families. These are existential, optimistic
statements; universal feasibility for every optimal response is different.
\end{proof}

\begin{proposition}[Sparse binary degrees and output length]
\label{prop:single-power-output}
Polynomial dependence on numerical degrees cannot in general be replaced by
polynomial dependence on their binary encoding, even with one power follower.
For an integer \(P\ge1\), \(z(x)=x^{1/P}\), \(0\le x\le1\), each of the following requires
\(\Omega(P)\) bits in any rational leader meeting the indicated guarantee:
\begin{enumerate}
\item upper objective \((z-1/2)^2\), no upper rows, additive error \(1/64\);
\item upper objective \(z\), upper row \(z\ge1/2\), exact feasibility and
additive error \(1/8\).
\end{enumerate}
The second example has a strict feasible anchor \(x=1\) and a convex reduced
upper inequality. Sparse upper data alone can also force exponential algebraic
output degree, with a trivial positive quadratic follower.
\end{proposition}
\begin{proof}
The power response minimizes \(z^{P+1}/(P+1)-xz\) on the unit interval,
with a sparse binary description of \(O(\log(P+1))\) bits.
In the first example the optimum is zero at \(x=2^{-P}\). The error guarantee
forces \(3/8\le z\le5/8\); hence \(0<x\le(5/8)^P\). In the second example
the optimum is \(1/2\), and the guarantee forces \(1/2\le z\le5/8\), giving
the same upper bound on positive \(x\). If \(x=p/q>0\) is a reduced rational,
then \(q\ge1/x\ge(8/5)^P\). Its binary encoding has \(\Omega(P)\) bits.
The reduced row \(1/2-x^{1/P}\le0\) is convex, and its value at one is
\(-1/2\). Thus an anchor does not remove the numerical-degree dependence of
Theorem~\ref{thm:accuracy-upper}.
For the last assertion let the upper feasible set be
\(1\le x\le2,\ x^{2^t}=2\), and minimize \(x\). Use the independent follower
\(\min_{0\le z\le1}z^2\), whose response is zero. The sparse upper description
has \(O(t)\) bits, whereas its unique leader has degree \(2^t\), since
\(X^{2^t}-2\) is Eisenstein at two.
\end{proof}
Example~\ref{ex:sparse-power-output} already gives the rational-output obstruction
with an affine upper objective and no upper rows by using two different powers.
Proposition~\ref{prop:single-power-output} strengthens the one-follower case
when polynomial upper data or a response row is allowed; it does not settle
the one-power affine-objective, no-upper-row subclass.

\begin{proposition}[An exponential common-field degree]\label{prop:root-sum-degree}
Independent uniformly strongly convex scalar polynomial followers can have an
aggregate value of algebraic degree \(2^N\). Even singleton convex quadratic
local sets can encode the exact Square Root Sum comparison problem through
one linear aggregate row.
\end{proposition}
\begin{proof}
For distinct primes \(p_i\), minimize \(z_i^3/3-p_i z_i\) on \([1,p_i+1]\).
The unique response is \(\sqrt{p_i}\), and the second derivative is at least
two throughout the interval. The value \(s=\sum_i\sqrt{p_i}\) has degree
\(2^N\). Here is a direct field argument. Inductively, the products
\(\sqrt{\prod_{i\in S}p_i}\) form a basis for the field generated by the first
\(j\) radicals, with independent sign-changing automorphisms. If the next
radical were in that field, its image under any such automorphism would be
itself or its negative, since its square is rational. It would be a common
eigenvector for all sign changes. In the displayed basis the characters are
distinct, so such an eigenvector is a rational multiple of one basis element.
Squaring contradicts the parity of the prime valuations of \(p_{j+1}\).
The next extension therefore has degree two and all old sign changes extend
with either new sign, completing the induction. The \(2^N\) signed sums are
distinct by this basis independence; they are all conjugates of \(s\).

For arbitrary nonnegative integers \(a_i\), the local set
\(Y_i=\{y_i:0\le y_i\le\max(1,a_i),\ y_i^2=a_i\}\) is a compact convex
singleton. Feasibility of \(\sum_i y_i\le K\) is exactly
\(\sum_i\sqrt{a_i}\le K\). There is no core variable and only one aggregate.
\end{proof}
The lower bound concerns minimal-polynomial and common-number-field output;
a list of radicals can be a much shorter expression. Exact Square Root Sum
comparison is not proved NP-hard by this encoding, and no such claim is used
here; polynomial-time exact comparison remains an open problem
\citep{EisenbrandEtAl2024}. It is an arithmetic obstruction to extending
Appendix~\ref{app:fixed-core}'s polyhedral-leaf elimination by replacing each
leaf with an arbitrary convex semialgebraic set. Accuracy-bit approximation,
which does not require the entire common field, remains compatible with it.
